\begin{theorem}[无外加寄存器与置换自然性约束下的二值跳跃结构唯一性]\label{thm:bdry-z2-jump-functor-uniqueness}
令 \(k\ge 2\)，并令 \(\mathbf{Set}_k^{\cong}\) 为由 \(k\) 元有限集及其双射构成的群胚。
令 \(\mathbf{Tors}(\ZZ_2)\) 为 \(\ZZ_2\)-torsor（带自由传递 \(\ZZ_2\) 作用的两点集合）及其同构构成的群胚。
称一个\emph{二值跳跃结构}为一个函子
\[
\mathcal J:\mathbf{Set}_k^{\cong}\longrightarrow \mathbf{Tors}(\ZZ_2).
\]
则当 \(k\ge 3\) 时，除处处平凡的函子外，存在且仅存在一个非平凡 \(\mathcal J\)（在自然同构意义下唯一）。
该唯一非平凡函子可由取取向构造给出：对任意 \(k\) 元集 \(S\)，定义
\[
\mathrm{Or}(S):=\mathrm{Bij}([k],S)\big/\mathfrak{A}_k,
\qquad [k]:=\{1,\dots,k\},
\]
其中 \(\mathfrak{A}_k\trianglelefteq\mathfrak{S}_k\) 为交错群，\(\mathfrak{S}_k\) 通过预合成作用于 \(\mathrm{Bij}([k],S)\)。
则 \(\mathfrak{S}_k\) 在 \(\mathrm{Or}(S)\) 上的作用恰通过符号同态 \(\operatorname{sgn}:\mathfrak{S}_k\to \ZZ_2\) 因子化，从而 \(\mathrm{Or}(S)\) 自然成为 \(\ZZ_2\)-torsor，且由 \(S\mapsto \mathrm{Or}(S)\) 定义的函子即为上述唯一非平凡二值跳跃结构。
\leanverified{signHom}
\leanverified{signHom\_exists}
\leanverified{trivHom}
\leanverified{sign\_range\_pm\_one}
\leanverified{sign\_sq\_eq\_one}
\leanverified{trivHom\_apply}
\leanverified{sign\_neq\_triv}
\leanverified{paper\_bdry\_z2\_jump\_functor\_skeleton}
\leanverified{paper\_bdry\_z2\_jump\_functor\_uniqueness}
\end{theorem}

\begin{proof}
取基对象 \(S_0=[k]\)，则 \(\mathrm{Aut}_{\mathbf{Set}_k^{\cong}}(S_0)\cong \mathfrak{S}_k\)。
任一 \(\ZZ_2\)-torsor 的自同构群同构于 \(\ZZ_2\)，因此对任意函子 \(\mathcal J\)，由 \(\mathfrak{S}_k\) 对 \(\mathcal J(S_0)\) 的自同构作用得到群同态
\[
\chi_{\mathcal J}:\mathfrak{S}_k\to \mathrm{Aut}(\mathcal J(S_0))\cong \ZZ_2.
\]
反之，给定任意群同态 \(\chi:\mathfrak{S}_k\to \ZZ_2\)，可令 \(\mathcal J_\chi(S_0)\) 为两点集合并以 \(\chi\) 指定其 \(\mathfrak{S}_k\)-作用；由于 \(\mathbf{Set}_k^{\cong}\) 为连通群胚，沿任意双射搬运可将该 \(\mathfrak{S}_k\)-作用推送到任意对象 \(S\)，从而得到一个函子 \(\mathcal J_\chi\)，并且该构造在自然同构意义下仅依赖 \(\chi\)。
因此二值跳跃结构的自然同构类与 \(\mathrm{Hom}(\mathfrak{S}_k,\ZZ_2)\) 一一对应。
\par
当 \(k\ge 3\) 时，\(\mathfrak{S}_k^{\mathrm{ab}}\cong \ZZ_2\)，故 \(\mathrm{Hom}(\mathfrak{S}_k,\ZZ_2)\) 仅含平凡同态与符号同态。
于是非平凡二值跳跃结构在自然同构意义下唯一。
最后，\(\mathrm{Or}(S)\) 的定义把 \(\mathrm{Bij}([k],S)\) 按 \(\mathfrak{A}_k\) 取商，得到一个两点集合；其上 \(\mathfrak{S}_k\) 的作用由 \(\mathfrak{S}_k/\mathfrak{A}_k\cong \ZZ_2\) 给出，正是符号同态所诱导的非平凡作用，从而对应唯一非平凡类。证毕。
\end{proof}

\begin{corollary}[boundary uplift 的取向 parity：从层标记到取向 torsor]\label{cor:bdry-uplift-orientation-parity}
在不引入外加寄存器且要求对 uplift 分支重标号自然的约束下，
对每个 boundary uplift 差集 \(\Delta_m^{\mathrm{bdry}}(w)\)（表 \ref{tab:foldbin_boundary_lift_m6_m8}），定义其取向 parity 为
\[
\Pi_m(w):=\mathrm{Or}\!\bigl(\Delta_m^{\mathrm{bdry}}(w)\bigr)\in \mathbf{Tors}(\ZZ_2).
\]
则：
\begin{enumerate}
  \item \(\Pi_m(w)\) 对 uplift 分支的任意重标号具有严格函子性；特别地，\(\Pi_m(w)\) 仅依赖 \(\Delta_m^{\mathrm{bdry}}(w)\) 作为无序有限集的同构类。
  \item 当 \(|\Delta_m^{\mathrm{bdry}}(w)|=2\) 时，\(\Pi_m(w)\) 退化为二点纤维上的唯一非平凡 \(\ZZ_2\)-torsor，并与 window-\(6\) 的二层 sheet parity（推论 \ref{cor:terminal-delta-parity-rule}）在自然同构意义下等价。
  \item 当 \(|\Delta_m^{\mathrm{bdry}}(w)|=3\) 时，\(\Pi_m(w)\) 给出无需指定特异层的唯一非平凡二值跳跃对象：它并非在“层集合本身”上的置换不变二分，而是由标号取向所诱导的 \(\ZZ_2\)-torsor，因而与命题 \ref{prop:bdry-sheet-parity-extension-obstruction-m7} 的不可延拓断言并不矛盾。
\end{enumerate}
\end{corollary}
\leanverified{upliftOrientationParity}
\leanverified{upliftOrientationParity\_relabel}
\leanverified{upliftOrientationParity\_two\_swap}
\leanverified{upliftOrientationParity\_three\_unique}
\leanverified{paper\_bdry\_uplift\_orientation\_parity}

\begin{proof}
第 (1) 条由 \(\mathrm{Or}:\mathbf{Set}_k^{\cong}\to\mathbf{Tors}(\ZZ_2)\) 的函子性直接给出。
第 (2) 条中，当 \(|\Delta|=2\) 时，\(\mathfrak{S}_2\cong\ZZ_2\) 的非平凡作用给出二点集上的唯一非平凡 torsor 结构；而 window-\(6\) 的 sheet parity 正是由二点纤维换位对合所诱导的同一 torsor 类型。
第 (3) 条中，\(|\Delta|=3\) 时由定理 \ref{thm:bdry-z2-jump-functor-uniqueness}（取 \(k=3\)）知非平凡二值跳跃结构在自然同构意义下唯一，并由 \(\mathrm{Or}\) 实现；同时命题 \ref{prop:bdry-sheet-parity-extension-obstruction-m7} 排除的是对三元集合\emph{置换不变的二分分级}（等价于选取特异层的 \(1\sqcup 2\) 分割），而 \(\mathrm{Or}(\Delta)\) 属于 torsor 型对象，二者类型不同。证毕。
\end{proof}

\begin{proposition}[uplift 纤维的符号降维：二值跳跃从中心迁移到交换化]\label{prop:uplift-fiber-sign-dimension-reduction}
\leanverified{paper\_gu\_uplift\_fiber\_sign\_dimension\_reduction}
令 \(S\) 为有限集，\(|S|=d\ge 2\)，并令
\[
\mathrm{Lab}(S):=\mathrm{Bij}([d],S)
\]
为其标号集合，则 \(\mathrm{Lab}(S)\) 是主 \(\mathfrak{S}_d\)-torsor。
记符号同态 \(\operatorname{sgn}:\mathfrak{S}_d\to \ZZ_2\)。
则存在自然同构
\[
\mathrm{Or}(S)\ \cong\ \mathrm{Lab}(S)\times_{\mathfrak{S}_d,\operatorname{sgn}}\ZZ_2.
\]
因此当 \(d=2\) 时，\(\mathfrak{S}_2\cong\ZZ_2\) 为阿贝尔群，二值跳跃可由纤维内换位对合实现为“中心型”翻转；
而当 \(d=3\) 时，\(Z(\mathfrak{S}_3)=\{e\}\) 但 \(\mathfrak{S}_3^{\mathrm{ab}}\cong\ZZ_2\) 仍由 \(\operatorname{sgn}\) 给出，从而二值跳跃只能以符号表示（交换化型）出现，即表现为取向 torsor，而非层标签。
\end{proposition}

\begin{proof}
由定义 \(\mathrm{Or}(S)=\mathrm{Bij}([d],S)/\mathfrak{A}_d\)。
另一方面，纤维积 \(\mathrm{Lab}(S)\times_{\mathfrak{S}_d,\operatorname{sgn}}\ZZ_2\) 等价于将两标号视为同类当且仅当它们相差一个偶置换；这正是对 \(\mathfrak{A}_d\) 的取商。
关于中心与交换化的群论断言为对称群的标准结构性质：当 \(d\ge 3\) 时 \(Z(\mathfrak{S}_d)=\{e\}\)，而当 \(d\ge 2\) 时 \(\mathfrak{S}_d^{\mathrm{ab}}\cong\ZZ_2\) 且由 \(\operatorname{sgn}\) 给出。证毕。
\end{proof}
